\subsection{Characterization of integral elements}

THEOREM 3. Let A be a subring of a field K. The integral closure A' of A in K is the intersection of the valuation rings of K which contain A; if A is local, A' is the intersection of the valuation rings of K which dominate A.

Let \(x\) be an element of \(A'\) and \(V\) a valuation ring of \(K\) containing \(A\) ; as \(x\) is integral over \(V\) , there exists a prime ideal \(p'\) of \(V[x]\) such that \(p' \cap V = m(V)\) (Chapter V, §2, no. 1, Theorem 1); clearly then the local ring \((V[x])_{p'}\) dominates \(V\) and hence is equal to it; whence \(x \in V\) . Conversely let \(y\) be an element of \(K\) which is not integral over \(A\) ; then there exists a maximal ideal \(M\) of \(A[y^{-1}]\) which contains \(y^{-1}\) (no. 2, Lemma 1); there exists also a valuation ring \(V\) of \(K\) which dominates \((A[y^{-1}])_{m}\) (no. 2, Corollary to Theorem 2); as \(y^{-1} \in m(V)\) , \(y \notin V\) . Further \(M \cap A\) is a maximal ideal of \(A\) (no. 2, Lemma 1); hence, if \(A\) is local, \(M \cap A = m(A)\) and \(V\) dominates \(A\) .

COROLLARY 1. Every valuation ring is integrally closed.

COROLLARY 2. For an integral domain to be integrally closed, it is necessary and sufficient that it be the intersection of a family of valuation rings of its field of fractions.

In the case of a Noetherian ring, Corollary 2 can be made more precise (Chapter VII, § 1, no. 3, Corollary to Theorem 1).

COROLLARY 3. Let K be a field, K' an extension of K and A a valuation ring of K. The integral closure of A in K' is the intersection of the valuation rings V' of K' such that V' ∩ K = A.

Theorem 1 (c) show that, if \(V'\) is a valuation ring of \(K'\) , \(V' \cap K\) is a valuation ring of K and \(V'\) dominates \(V' \cap K\) . For \(V'\) to dominate A, it is necessary and sufficient that \(V' \cap K\) dominate A and therefore be equal to it.

